% Source: https://artofproblemsolving.com/community/c3h3083065_absolute_values_with_variable_powers_inside_radicals
% Stable identity: aops-topic-3083065-q01
At $ x=3 $, both powers have positive exponents and both sides are zero, so it is a solution. Otherwise the base $ b=|x-3|$ is positive, and the equation is $ b^{\dfrac{x+1}{4}}=b^{\dfrac{x-2}{3}}$. If $ b=1 $, then $ x=2 $ or $ 4 $, both solutions. If $ b\ne1 $, exponential injectivity gives $\dfrac{x+1}{4}=\dfrac{x-2}{3}$, hence $ x=11 $. It has base $ 8 $ and satisfies the equation. The complete set is $\boxed{\{2,3,4,11\}}$.
